\section{Related Work}
\subsection{Techniques for Large Language Model Compression}
The growing size of large language models (LLMs) has raised increasing concerns over their computational and memory demands. 
To address these challenges, a variety of model compression techniques have been proposed. 
Conventional approaches often require computationally expensive retraining, which is generally impractical due to the substantial computational cost associated with the massive size of LLMs. 
Consequently, recent efforts have shifted toward more resource-friendly post-training compression techniques~\cite{frantar2022gptq,frantar2023sparsegpt,zhu2024survey}. 
Typically employed approaches include network pruning and weight quantization. And pruning techniques can be classified into unstructured and structured methods.
Unstructured pruning removes individual weights based on importance scores. 
For example, SparseGPT~\cite{frantar2023sparsegpt} performs one-shot pruning using second-order approximations without retraining.  
However, since unstructured pruning retains the original matrix shape, it offers limited inference acceleration and requires specialized hardware. 
In contrast, structured pruning eliminates entire blocks or channels, enabling compatibility with conventional hardware platforms. 
LLM-Pruner~\cite{ma2023llm} groups dependent linear projections into coupled structures, assigns each group a loss‑aware importance score, prunes the least important groups, and applies LoRA fine‑tuning to restore performance. Additionally, ZipLM~\cite{kurtic2023ziplm} prioritizes pruning components that yield the worst trade-off between latency and loss, but this often causes notable performance degradation. 
Quantization provides another mainstream solution. 
GPTQ~\cite{frantar2022gptq} gradually quantizes and updates each weight using the Hessian matrix to minimize the quantization error. 
AWQ~\cite{lin2024awq} preserves important weight channels by selecting reparameterization coefficients via grid search. Furthermore, BiLLM~\cite{huang2024billm} and ARB-LLM~\cite{li2024arb} push quantization to the 1-bit level, while still delivering impressive downstream task accuracy. 
However, these techniques still cause significant performance degradation, especially at low bit-widths, due to the lack of weight and activation adaptation. 

\subsection{SVD-based Techniques for Compressing LLMs}
Singular value decomposition (SVD) is a commonly employed technique for reducing matrix dimensionality by representing the original matrix as the product of two low-rank factor matrices. 
Recent research has demonstrated the potential of SVD-based LLM compression methods, yet comprehensive exploration remains limited. 
Hsu et al.~\cite{hsulanguage} propose FWSVD, incorporating Fisher information to weight the importance of parameters. 
However, it relies on computationally intensive training and was originally applied only to small language models (e.g., BERT~\cite{devlin2019bert}, ALBERT~\cite{lan2019albert}). 
LoRD~\cite{kaushal2023lord} first applies SVD to LLMs by grouping layers to improve efficiency. 
However, LoRD overlooks the importance of input activations.
To address this limitation, Yuan et al.~\cite{yuan2023asvd} introduce ASVD, which mitigates the impact of activation outliers by reshaping the weight matrix based on activation distribution, thereby enhancing the precision of the low-rank decomposition. 
CALDERA~\cite{saha2024compressing} exploits the inherent low-rank structure of weight matrices by approximating them through a low-rank, low-precision decomposition.
Additionally, SVD-LLM~\cite{wang2024svd} establishes a direct mapping between singular values and truncation loss, which means that truncating the smallest singular values leads to minimal truncation loss.
Basis Sharing~\cite{wang2024basis} achieves more effective LLM compression through sharing parameters across different layers. 
More recently, AdaSVD~\cite{li2025adasvd} proposes an adaptive compensation method, iteratively updating matrices during truncation. 

Despite these advances, all existing methods ignore the residual matrix produced during SVD truncation, which can significantly compensate for the SVD truncation loss. 
Moreover, compressing only the last few layers of the model under a fixed target compression ratio can provide better performance for compressed models. 
However, existing methods compress all layers, either apply a uniform compression ratio across all layers or assign variable ratios based on layer-wise importance, which often results in sub-optimal performance.
